\begin{helper}
Let \(D\) be a digraph on a finite vertex set \(W\), and let \(3\le \ell\).
Put
\[
  \Omega=\{\phi:\ [\ell-1]\to([n]\to W)\}.
\]
Call \(\phi\in\Omega\) bad if the coloring
\(\noderef{RandomHomomorphismColoring}(D,\ell,n,\phi)\) has a final-color
monochromatic clique on some \(s\)-element set.  Then the number of bad maps is
at most
\[
  \binom ns\,
  \overrightarrow{i_s}(D)^{\ell-1}\,
  |W|^{(\ell-1)(n-s)}.
\]
\end{helper}
\begin{proof}
For each bad \(\phi\), choose one witnessing \(s\)-element set \(S\).  This
embeds the bad maps into the set of pairs \((\phi,S)\) where \(|S|=s\) and
every edge inside \(S\) has final color under
\(\noderef{RandomHomomorphismColoring}\).

Fix such a pair \((\phi,S)\), and list \(S\) increasingly as
\[
  v_1<\cdots <v_s .
\]
For every first color \(c<\ell\), the final-color assumption on all edges of
\(S\) implies, by \(\noderef{RandomHomomorphismFinalColorForwardIndependent}\),
that the tuple
\[
  (\phi_c(v_1),\ldots,\phi_c(v_s))
\]
is forward independent in \(D\).  Thus the pair \((\phi,S)\) determines:
the set \(S\); for each \(c<\ell\), one forward-independent \(s\)-tuple in
\(D\); and for each \(c<\ell\), all values of \(\phi_c\) outside \(S\).

This data determines \(\phi\) uniquely.  Indeed, on \(S\) the increasing list
identifies every point of \(S\) with a unique index \(i\in[s]\), so the chosen
forward-independent tuple gives \(\phi_c(v_i)\).  Outside \(S\), the outside
values are recorded directly.  Hence the number of such pairs is bounded by
\[
  \#\{S\subseteq[n]: |S|=s\}\,
  \overrightarrow{i_s}(D)^{\ell-1}\,
  |W|^{(\ell-1)(n-s)}.
\]
There are exactly \(\binom ns\) choices of \(S\), and
\(\noderef{ForwardIndependentTupleCount}\) identifies
\(\overrightarrow{i_s}(D)\) with the number of forward-independent
\(s\)-tuples.  This gives the claimed upper bound.
\end{proof}

